\documentclass[12pt,a4paper]{article}


\usepackage[utf8]{inputenc}
\usepackage[T1]{fontenc}
\usepackage[polish]{babel}
\usepackage{lmodern}


\usepackage[margin=2.5cm]{geometry}
\usepackage{setspace}
\onehalfspacing
\usepackage{indentfirst}         


\usepackage{fancyhdr}
\pagestyle{fancy}
\fancyhf{}
\lhead{Bazy danych I}
\rhead{System obsługi hotelu}
\cfoot{\thepage}


\usepackage{booktabs}            
\usepackage{tabularx}             
\usepackage{longtable}            
\usepackage{array}
\usepackage{multirow}


\usepackage{graphicx}
\usepackage{float}              
\usepackage{caption}
\usepackage{subcaption}
\graphicspath{{img/}}


\usepackage{enumitem}


\usepackage{xcolor}
\usepackage{listings}
\lstset{
  language=SQL,
  basicstyle=\ttfamily\small,
  keywordstyle=\color{blue}\bfseries,
  commentstyle=\color{gray}\itshape,
  stringstyle=\color{teal},
  numbers=left,
  numberstyle=\tiny\color{gray},
  frame=single,
  breaklines=true,
  showstringspaces=false,
  tabsize=2,
  morekeywords={EXCLUDE, USING, GIST, RETURNS, TRIGGER, LANGUAGE, PLPGSQL,
                RETURN, NEW, OLD, BEGIN, END, IF, THEN, RAISE, EXCEPTION,
                daterange, numeric, serial},
  literate={ą}{{\k{a}}}1 {ę}{{\k{e}}}1 {ć}{{\'c}}1 {ł}{{\l{}}}1
           {ń}{{\'n}}1 {ó}{{\'o}}1 {ś}{{\'s}}1 {ź}{{\'z}}1 {ż}{{\.z}}1
           {Ą}{{\k{A}}}1 {Ę}{{\k{E}}}1 {Ć}{{\'C}}1 {Ł}{{\L{}}}1
           {Ń}{{\'N}}1 {Ó}{{\'O}}1 {Ś}{{\'S}}1 {Ź}{{\'Z}}1 {Ż}{{\.Z}}1
}
\renewcommand{\lstlistingname}{Listing}


\usepackage[hidelinks]{hyperref}

\begin{document}
\begin{titlepage}
  \centering
  \vspace*{\fill}

  {\Large Bazy danych I -- projekt zaliczeniowy\par}
  \vspace{0.8cm}
  {\Huge\bfseries HotelLukas\par}
  \vspace{0.5cm}
  {\large System obsługi hotelu -- dokumentacja projektu\par}

  \vspace{2.5cm}

  {\large Autor: \textbf{Łukasz Deja}\par}

  \vspace{1cm}

  {\normalsize Repozytorium projektu:\par}
  \vspace{0.2cm}
  {\normalsize\url{https://github.com/lukaszdeja/HotelDataBase}\par}

  \vspace*{\fill}
\end{titlepage}

\section*{Część I: Projekt koncepcji i założenia}
\addcontentsline{toc}{section}{Część I: Projekt koncepcji i założenia}

\section{Temat projektu}

\noindent Projekt realizuje system zarządzania hotelem. Jego celem jest umożliwienie rezerwacji pokoi hotelowych przez gości, wykrywanie kolizji przy rezerwacjach. Dodatkowo umożliwia generowanie i przeglądanie raportów z przychodów dla administratorów. W wersji rozszerzonej umożliwia także rezerwacje usług dodatkowych. Relacyjna baza danych PostgreSQL a także aplikacja webowa obsługują podział na role użytkowników.

Zadania systemu:
\begin{itemize}
    \item Rejestracja konta i logowanie.
    \item Rezerwacje pokojów oraz przechowywanie w bazie danych informacji o gościach oraz pokojach.
    \item Rejestrowanie płatności (bez faktycznego transferu środków, tylko wpis w bazie).
    \item Generowanie raportów - obłożenia pokoi, miesięcznych przychodów
    \item Podział na role - gość, recepcjonista i administrator
    \item \textbf{W wersji rozszerzonej: } system usług hotelowych (sauna, basen), rezerwacje usług
\end{itemize}

\section{Analiza wymagań użytkownika}

\subsection{Historyjki użytkownika i priorytety}

\noindent Wymagania zebrano w postaci historyjek użytkownika (backlog) dla trzech ról:
gościa, recepcjonisty i administratora. Każdej historyjce nadano priorytet
według metody MoSCoW: M (\emph{must} -- konieczne), S (\emph{should} --
ważne), C (\emph{could} -- opcjonalne), W (\emph{won't} -- poza wersją
podstawową; historyjki W realizuje wersja rozszerzona). Ostatnia kolumna
wskazuje wymaganie funkcjonalne (punkt \ref{sec:wymagania}), które realizuje
daną historyjkę.

\begin{longtable}{>{\raggedright\arraybackslash}p{0.7cm}
                  >{\raggedright\arraybackslash}p{11cm}
                  >{\centering\arraybackslash}p{0.9cm}
                  >{\raggedright\arraybackslash}p{1.6cm}}
\toprule
\textbf{Nr} & \textbf{Historyjka} & \textbf{Prio.} & \textbf{Wym.} \\
\midrule
\endhead

\multicolumn{4}{l}{\textit{Gość}} \\
\midrule
1  & Jako gość chcę utworzyć konto, podając swoje dane kontaktowe, aby móc rezerwować pokoje.
   & M & \ref{f:rejestracja} \\
2  & Jako gość chcę zalogować się i wylogować, aby móc korzystać z serwisu.
   & M & \ref{f:logowanie} \\
3  & Jako gość chcę zmienić swoje dane kontaktowe, aby hotel mógł się ze mną skontaktować.
   & C & \ref{f:zmianadanych} \\
4  & Jako gość chcę wyszukać wolne pokoje w wybranym terminie, aby wiedzieć, co i kiedy mogę zarezerwować.
   & M & \ref{f:wolne} \\
5  & Jako gość chcę zawęzić wyszukiwanie do wybranego typu pokoju i liczby osób, aby znaleźć pokój dopasowany do moich wymagań.
   & S & \ref{f:wolne} \\
6  & Jako gość chcę zarezerwować jeden lub więcej pokoi, aby móc spędzić wakacje w hotelu.
   & M & \ref{f:rezerwacja} \\
7  & Jako gość chcę mieć pewność, że system gwarantuje brak kolizji terminów, aby nikt nie zarezerwował mojego pokoju w moim terminie.
   & M & \ref{f:kolizje} \\
8  & Jako gość chcę poznać koszt pobytu przed potwierdzeniem rezerwacji, aby wiedzieć, ile zapłacę.
   & S & \ref{f:koszt} \\
9  & Jako gość chcę przeglądać swoje rezerwacje z sumą wpłat i statusem płatności, aby wiedzieć, ile zostało do zapłaty.
   & M & \ref{f:status} \\
10 & Jako gość chcę anulować rezerwację przed przyjazdem, jeśli jeszcze nic nie wpłaciłem, aby nie płacić za pobyt, z którego rezygnuję.
   & S & \ref{f:anulowanie} \\
11 & Jako gość chcę zarezerwować usługę (saunę, basen) na konkretny dzień i godzinę w trakcie pobytu, aby nie czekać w kolejce.
   & W & \ref{f:rezerwacjauslugi} \\
12 & Jako gość chcę widzieć łączny koszt pobytu razem z usługami, aby znać pełną kwotę do zapłaty.
   & W & \ref{f:kosztuslug} \\
\midrule
\multicolumn{4}{l}{\textit{Recepcjonista}} \\
\midrule
13 & Jako recepcjonista chcę zalogować się do systemu na konto recepcjonisty, aby móc pracować.
   & S & \ref{f:logowanierecepcja} \\
14 & Jako recepcjonista chcę przeglądać i edytować dane gości, aby kontrolować, czy są wprowadzone poprawnie.
   & C & \ref{f:danegosci} \\
15 & Jako recepcjonista chcę przeglądać wszystkie rezerwacje, aby kontrolować sytuację hotelu.
   & M & \ref{f:danegosci} \\
16 & Jako recepcjonista chcę zaksięgować zaliczkę i dopłatę za rezerwację, aby wiedzieć, które rezerwacje są już opłacone.
   & M & \ref{f:wplaty} \\
17 & Jako recepcjonista chcę, aby nie dało się wpłacić więcej niż pozostały koszt, aby uniknąć nadpłat i zwrotów.
   & S & \ref{f:nadplata} \\
\midrule
\multicolumn{4}{l}{\textit{Administrator}} \\
\midrule
18 & Jako administrator chcę zalogować się do systemu na konto administratora, aby zarządzać hotelem.
   & S & \ref{f:logowanieadmin} \\
19 & Jako administrator chcę dodawać, zmieniać i usuwać pokoje, aby oferta odpowiadała stanowi hotelu.
   & S & \ref{f:pokoje} \\
20 & Jako administrator chcę zarządzać typami pokoi i ich cenami, aby móc dostosowywać cennik.
   & S & \ref{f:pokoje} \\
21 & Jako administrator chcę zarządzać listą usług hotelowych i ich cenami, aby dostosowywać ofertę hotelu.
   & W & \ref{f:uslugi} \\
22 & Jako administrator chcę przeglądać raport przychodów w poszczególnych miesiącach, aby śledzić wyniki finansowe hotelu.
   & S & \ref{f:przychody} \\
23 & Jako administrator chcę przeglądać raport obłożenia według typu pokoju, aby wiedzieć, które pokoje cieszą się największym zainteresowaniem.
   & S & \ref{f:oblozenie} \\
24 & Jako administrator chcę przeglądać listę gości o najwyższych wydatkach, aby móc proponować im rabaty.
   & C & \ref{f:najlepsigoscie} \\
25 & Jako administrator chcę przeglądać listę rzadko rezerwowanych pokoi, aby zdecydować o zmianie ich ceny.
   & C & \ref{f:rzadkie} \\
26 & Jako administrator chcę wiedzieć, w których miesiącach przychód przekroczył średnią, aby planować promocje na słabsze miesiące.
   & C & \ref{f:ponadsrednia} \\
27 & Jako administrator chcę tworzyć i usuwać konta recepcjonistów, aby nowi pracownicy mieli dostęp do systemu.
   & S & \ref{f:kontapracownikow} \\
\bottomrule
\end{longtable}

Minimalny zakres działającego systemu (MVP) tworzą historyjki o priorytecie M:
1, 2, 4, 6, 7, 9, 15 i 16. Pierwsze konto administratora tworzone jest
w skrypcie z danymi początkowymi; kolejne konta pracowników zakłada
administrator z poziomu aplikacji.

\subsection{Wymagania funkcjonalne -- wersja podstawowa}
\label{sec:wymagania}

\subsubsection*{Gość}
\begin{enumerate}[label=F\arabic*., ref=F\arabic*, series=wymagania]
    \item Gość może utworzyć konto, podając imię, nazwisko, e-mail, telefon i hasło. \label{f:rejestracja}
    \item Gość może zalogować się do systemu i wylogować z niego. \label{f:logowanie}
    \item Gość może zmienić swoje dane kontaktowe. \label{f:zmianadanych}
    \item Gość może wyszukiwać wolne pokoje w wybranym przez siebie terminie, z możliwością
          filtrowania po typie pokoju i liczbie osób. \label{f:wolne}
    \item Gość może zarezerwować jeden lub kilka wolnych pokoi na wybrany
          termin. \label{f:rezerwacja}
    \item Gość może przeglądać swoje rezerwacje wraz z ich kosztem,
          sumą dokonanych wpłat i statusem płatności (brak wpłat,
          opłacona częściowo, opłacona w całości). \label{f:status}
    \item Gość może anulować swoją rezerwację przed datą rozpoczęcia pobytu o ile nie dokonano jeszcze żadnej wpłaty. \label{f:anulowanie}
\end{enumerate}

\subsubsection*{Recepcjonista}
\begin{enumerate}[label=F\arabic*., ref=F\arabic*, resume=wymagania]
    \item Recepcjonista może zalogować się do systemu i wylogować z niego. \label{f:logowanierecepcja}
    \item Recepcjonista może przeglądać i edytować dane gości oraz
          wszystkie rezerwacje. \label{f:danegosci}
    \item Recepcjonista może księgować wpłaty za rezerwacje -- zaliczkę oraz dopłatę pozostałej należności. \label{f:wplaty}
\end{enumerate}

\subsubsection*{Administrator}
\begin{enumerate}[label=F\arabic*., ref=F\arabic*, resume=wymagania]
    \item Administrator może zalogować się do systemu i wylogować z niego. \label{f:logowanieadmin}
    \item Administrator może tworzyć, edytować i usuwać konta recepcjonistów. \label{f:kontapracownikow}
    \item Administrator może dodawać, edytować i usuwać pokoje oraz typy pokoi wraz z ich cenami. \label{f:pokoje}
    \item Administrator może generować raport miesięcznych przychodów. \label{f:przychody}
    \item Administrator może generować raport obłożenia pokoi według typu. \label{f:oblozenie}
    \item Administrator może generować raport gości o najwyższych
          wydatkach. \label{f:najlepsigoscie}

    \item Administrator może generować raport pokoi rzadko rezerwowanych. \label{f:rzadkie}
    \item Administrator może generować raport miesięcy, w których przychód przekroczył średnią miesięczną. \label{f:ponadsrednia}
\end{enumerate}

\subsubsection*{System}
\begin{enumerate}[label=F\arabic*., ref=F\arabic*, resume=wymagania]
    \item System uniemożliwia rezerwację pokoju zajętego w danym terminie. \label{f:kolizje}
    \item System automatycznie wylicza koszt rezerwacji na podstawie liczby nocy i ceny pokoju w momencie rezerwacji. \label{f:koszt}
    \item System odrzuca niepoprawne dane (np. data wymeldowania wcześniejsza niż data zameldowania, ujemna kwota płatności). \label{f:walidacja}
    \item System uniemożliwia zaksięgowanie wpłaty przekraczającej
    pozostałą należność za rezerwację. \label{f:nadplata}
\end{enumerate}

\subsection{Wymagania funkcjonalne -- wersja rozszerzona}
\begin{enumerate}[label=F\arabic*., ref=F\arabic*, resume=wymagania]
    \item Administrator może zarządzać listą usług hotelowych i ich cenami. \label{f:uslugi}
    \item Gość może zarezerwować usługę na konkretny dzień i godzinę w trakcie swojego pobytu.. \label{f:rezerwacjauslugi}
    \item Koszt zarezerwowanych usług jest doliczany do kosztu pobytu. \label{f:kosztuslug}
\end{enumerate}

\subsection{Wymagania niefunkcjonalne}
\begin{itemize}
    \item Technologie: baza danych PostgreSQL, aplikacja webowa (Node.js)
    \item Bezpieczeństwo: podział na role, również na poziomie bazy danych, hashowanie haseł, tokeny logowania.
    \item Automatyzacja: triggery w bazie danych.
    \item Szybkość i wydajność: System sprawnie obsługuje duże zapytania dzięki indeksom.
    \item Dostępność: aplikacja webowa ORAZ baza danych (obiekty w bazie) w \textbf{języku polskim.}
\end{itemize}

\section{Funkcje realizowane w bazie danych}

\noindent Poniższa tabela przedstawia funkcje realizowane po stronie bazy danych.


\begin{longtable}{>{\raggedright\arraybackslash}p{4.5cm}
                  >{\raggedright\arraybackslash}p{3cm}
                  >{\raggedright\arraybackslash}p{5.5cm}
                  p{1.5cm}}
\toprule
\textbf{Funkcja} & \textbf{Mechanizm} & \textbf{Opis} & \textbf{Wym.} \\
\midrule
\endhead

\multicolumn{4}{l}{\textit{Wersja podstawowa}} \\
\midrule
Wykrywanie kolizji rezerwacji & ograniczenie \texttt{EXCLUDE}
  & odrzuca rezerwację pokoju w terminie nakładającym się na inną
  & \ref{f:kolizje} \\
Wyszukiwanie wolnych pokoi & funkcja
  & zwraca pokoje wolne w podanym terminie, z filtrem typu i liczby osób
  & \ref{f:wolne} \\
   Wolne pokoje na dziś & widok, podzapytanie
     & pokoje wolne w bieżącym dniu z typem, liczbą osób i ceną,
       bez danych osób rezerwujących i płatności
     & \ref{f:wolne} \\
Wyliczanie kosztu rezerwacji & wyzwalacz, funkcja
  & wyzwalacz kopiuje aktualną cenę pokoju do pozycji rezerwacji;
    funkcja \texttt{koszt\_rezerwacji} liczy koszt pobytu na podstawie
    liczby nocy i zapisanej ceny, wykorzystywana w widoku statusu
    płatności i w blokadzie nadpłaty
  & \ref{f:koszt} \\
Walidacja danych & ograniczenia \texttt{CHECK}
  & poprawność dat, kwot, cen, formatu e-mail i telefonu, rodzaju
    płatności i roli pracownika
  & \ref{f:walidacja} \\
Blokada nadpłaty & wyzwalacz
  & odrzuca wpłatę przekraczającą pozostałą należność
  & \ref{f:nadplata} \\
Blokada anulowania opłaconej rezerwacji & klucz obcy (\texttt{RESTRICT})
  & uniemożliwia usunięcie rezerwacji, do której dokonano wpłaty
  & \ref{f:anulowanie} \\
  Blokada anulowania rozpoczętego pobytu & wyzwalacz
  & uniemożliwia usunięcie rezerwacji, której pobyt już się rozpoczął
  & \ref{f:anulowanie} \\
Usuwanie pozycji anulowanej rezerwacji & klucz obcy (\texttt{CASCADE})
  & usunięcie rezerwacji usuwa przypisane do niej pokoje
  & \ref{f:anulowanie} \\
Status płatności & widok, grupowanie
  & koszt, suma wpłat, kwota do zapłaty i status (brak, częściowo,
    w całości) dla każdej rezerwacji
  & \ref{f:status}, \ref{f:danegosci} \\
Raport przychodów & widok, grupowanie
  & przychody zgrupowane według miesięcy
  & \ref{f:przychody} \\
Raport obłożenia & widok, grupowanie
  & liczba zajętych nocy i procent obłożenia według typu pokoju
  & \ref{f:oblozenie} \\
Raport najlepszych gości & widok, grupowanie
  & goście o łącznych wydatkach powyżej progu
    (\texttt{HAVING} na \texttt{SUM})
  & \ref{f:najlepsigoscie} \\
Raport pokoi rzadko rezerwowanych & widok, grupowanie
  & pokoje z liczbą rezerwacji poniżej progu
    (\texttt{HAVING} na \texttt{COUNT})
  & \ref{f:rzadkie} \\
Raport miesięcy ponad średnią & widok, grupowanie, podzapytanie
  & miesiące o przychodzie wyższym od średniej miesięcznej
    (\texttt{HAVING} z podzapytaniem)
  & \ref{f:ponadsrednia} \\
Poziomy dostępu & widoki i uprawnienia
  & osobne widoki dla gościa, recepcjonisty i administratora,
    udostępniane poleceniem \texttt{GRANT}; prawo do widoku
    nie daje dostępu do tabel bazowych
  & \ref{f:rejestracja}--\ref{f:ponadsrednia} \\

\midrule
\multicolumn{4}{l}{\textit{Wersja rozszerzona}} \\
\midrule
Rezerwacja usługi w trakcie pobytu & wyzwalacz
  & odrzuca rezerwację usługi poza terminem pobytu gościa
  & \ref{f:rezerwacjauslugi} \\
Doliczanie kosztu usług & wyzwalacz, widok
  & wyzwalacz kopiuje aktualną cenę usługi do jej rezerwacji;
    widok statusu płatności dolicza koszt usług do kosztu pobytu
  & \ref{f:kosztuslug} \\
\bottomrule
\end{longtable}

\section*{Część II: Projekt diagramów}
\addcontentsline{toc}{section}{Część II: Projekt diagramów}

\end{document}